\scanpage{253}
\mspage{427}

This equation has the solution
\[
(63)\qquad\xi=\ell_1^\nu
\]
in $\widehat\Pi_{0,3}$ (for $\tau_\infty(\ell_1)=\ell_1^{-1}$
hence $\tau_\infty(\xi)=\ell_1^{-\nu}$)
and this solution is unique \uline{\unsure{if|once} we admit}
that the discrete results extend to the profinite case,
so as to have
\[
(64)\qquad\bigl(\widehat\Pi_{0,3}\bigr)^{\tau_\infty}=\{1\}.
\]
Now $\xi$ must moreover satisfy the relation
\[
\xi\rho(\xi)
\rho^2(\xi)=1
\]
and we have just seen that this implies
$\nu=\pm1$, i.e.
$\unsure{\dot u|u}\in\{\unsure{\mathrm{id}|1},\tau_\infty\}$.
\mspage{252}

\scanpage{254}
\mspage{428}

48. Construction of a group scheme.

I) Study of the map $\alpha\longmapsto
\alpha\rho(\alpha)^{-1}=[\alpha,\rho]$.

\unsure{Let|Let} $G$ be a group scheme over a ring
$A$ (for example $A=\mathbb Z$, $G=\mathrm{Sl}(2)_{\inserted{\mathbb Z}}
\deleted{\unsure{,\mathbb Z)}}$ or $\mathrm{Gl}(2)_{\mathbb Z}$),
$\rho$ a section of $G$, and let us consider the morphism
\[
(1)\qquad
F_\rho:\ \alpha\longmapsto\alpha\rho(\alpha)^{-1}
=[\alpha,\rho]:\qquad G\longrightarrow G.
\]
Let
\[
\deleted{\unsure{Z_\rho}}\quad
T_\rho
=\mathrm{Centr}_G(\rho)=G^\rho\subset G
\]
then
\[
(2)\qquad
F_\rho(\alpha')=F_\rho(\alpha)\quad\Longleftrightarrow\quad
\deleted{\unsure{\alpha}}\alpha'\deleted{\text{is}}
\text{ section of }\alpha T_\rho
\]
\inserted{thus} we have a canonical monomorphism
\inserted{induced by $F_\rho$}
\[
(3)\qquad
G/T_\rho\hookrightarrow G,
\]
whose image is to be characterized.

In the case that concerns us, the quotient group $G_{\mathrm{ab}}$
is representable (it is the trivial group for
$G=\mathrm{Sl}(2)_{\mathbb Z}$, and it is $\mathbb G_m$ if
$G=\mathrm{Gl}(2)_{\mathbb Z}$),
\deleted{and}\inserted{In any case,} a first condition for
a section $\xi$\inserted{$\in G(S)$}
over an $A$Scheme $S$, \deleted{over $A$} for it
to be in the image, is that we have
\[
(4)\qquad\delta(\xi)=1
\]
\mspage{253}

\scanpage{255}
\[
(5)\qquad \delta:G\longrightarrow G_{\mathrm{ab}}
\]
is the canonical \unsure{homomorphism}. A second condition is that $\xi\rho$, which must be (locally) of the form $\alpha\rho\alpha^{-1}$, be \uline{conjugate} \uline{locally} to $\rho$ -- and this condition is moreover necessary and sufficient, tautologically, for $\xi$ to belong to the image of $F_\rho$, and thus implies condition (4). But the latter may sometimes serve as an ingredient, \deleted{this} in a \illegible{1} of conditions of a ``computational'' nature, which are necessary and sufficient for us to have $\xi\rho\sim\rho$.

In the case of a reductive group $G$, the condition of \deleted{the \unsure{equivalence}}\inserted{\unsure{conjugation}} \inserted{locally} of sections is moreover fairly well \unsure{understood}, since in particular the work of Steinberg. We have a ``scheme of invariants'' $I_G$, and a canonical morphism
\[
(6)\qquad G\longrightarrow I_G
\]
\unsure{(}which is universal \inserted{(over $A$, and also after every base change)} for morphisms from $G$ into schemes, which are \unsure{invariant} under inner \unsure{automorphisms} -- i.e. $I_G$ is the \uline{representable} \uline{envelope} of $G/\mathrm{int}(G)$ \unsure{sheaf} of conjugacy
\mspage{254}

\scanpage{256}
\mspage{429}
classes.
\[
(7)\qquad G/\mathrm{int}(G)\longrightarrow I_G.
\]
To tell the truth, \deleted{it \illegible{1}} this morphism is an epimorphism (for \unsure{the|the} fppf topology), but not an isomorphism i.e. not a monomorphism -- in other words, for two sections (in this case $\rho$, $\xi\rho$) of $G$, it is indeed necessary, but not always sufficient, that their images in $I_G$ be equal, for them to be conjugate locally$^{t}$ fppf. But in the case where the two sections are \uline{\deleted{quasi}-regular} in the sense of Steinberg (which is a condition on each fiber) \deleted{it is}\inserted{the condition is} \unsure{sufficient}. It is also so if the two sections are \unsure{$\frac{1}{2}$} simple (which is also a condition on each fiber).

In the case of a group locally isomorphic to $\mathrm{Sl}(2)$, \inserted{to $\mathrm{GP}(1)$,} or $\mathrm{Gl}(2)$, the \uline{regular} sections in the sense of Steinberg are moreover those which are central at no point -- which is indeed the case for the $\rho$ or the $\sigma_\infty=\sigma$ of the preceding \S\. On the other hand, $I_G$ is then well known, \deleted{in the case of a group}
\mspage{255}

\scanpage{257}
\deleted{split, more generally of \uline{$\mathrm{Sl}(E)$} or \uline{$\mathrm{Gl}(E)$},}
$E$ locally free of rank 2, it is canonically isomorphic to $\mathbb E^1$ or to $\mathbb E^1\times\mathbb G_m$ (depending on whether the rank is 1 or 2) in the case of \deleted{\unsure{$\mathrm{Sl}(2)$|$\mathrm{Gl}(2)$}}a form of $\mathrm{Sl}(2)$ \deleted{\unsure{or $\mathrm{GP}(1)$}}, it is given by the \uline{\unsure{Trace}}
\[
(7)\qquad
\begin{array}{rcl}
G&\overset{\deleted{\mathrm{Tr}\;\unsure{\det}}}{\longrightarrow}&I_G\simeq\mathbb E^1,\\
x&\longmapsto&\mathrm{Tr}\,x
\end{array}
\]
in the case of a form of $\mathrm{Gl}(2)$ it is given by
\[
(8)\qquad
\begin{array}{rcl}
G&\longrightarrow&I_G=\mathbb E^{1*}\times\mathbb E^1.\\
x&\longmapsto&(\det x,\mathrm{tr}\,x)
\end{array}
\]
We thus find (for the record)

\uline{Proposition} Let $G$ be a form of $\mathrm{Sl}(2)$ \unsure{or} $\mathrm{Gl}(2)$ over a scheme $S$, let $\rho$ be a section of $G$ which is central in no fiber. For a section $\xi$ of $G$ to be \inserted{(locally)} in the image of $F_\rho$ ((1) or (3)), it is necessary and sufficient that it satisfy the following conditions

a) $\deleted{\unsure{\rho}}\xi\rho$ is central in no fiber, i.e. in no fiber do we have $\xi=\lambda\rho^{-1}$, where $\lambda$ \deleted{\unsure{is|being}} a central element

b) $\mathrm{tr}\,\xi\rho=\mathrm{tr}\,\rho$

c) (if $G$ a form of $\mathrm{Gl}(2)$) $\det\xi=1$ i.e. $\xi$ a section of $\mathrm{Sl}(2)$.

\margin{is equivalent to: $\det(\xi\rho)=\det\rho$}
\mspage{256}

\scanpage{258}
\mspage{430}
\uline{Corollary} When these conditions are satisfied, the scheme of solutions \inserted{in $\alpha$} of the equation
\[
[\alpha,\rho]=\xi
\]
is a (right) torsor under $T_\rho$, which is a smooth commutative group scheme over $S$, of relative dimension 1 (\unsure{if} $G$ a form of $\mathrm{Sl}(2)$) or 2 (if $G$ a form of $\unsure{\mathrm{Gl}(2)|\mathrm{Sl}(2)}$).

\uline{Examples \deleted{\unsure{1}}} $G=\mathrm{Gl}(2)_{\mathbb Z}$, $\rho=\begin{pmatrix}0&1\\-1&1\end{pmatrix}$ is the element of $\mathrm{Sl}(2,\mathbb Z)$ already so denoted (starting from page~335). It is indeed regular, \deleted{\unsure{and}} likewise the element $\sigma=\sigma_\infty=\begin{pmatrix}0&-1\\1&0\end{pmatrix}$.

In general, the datum of a form of $\mathrm{Gl}(2)$ is equivalent to that of a form of $\mathrm{Sl}(2)$, or of $\mathrm{GP}(1)$, or again to that of an Azumaya Algebra of rank 4 over $S$, let $M$ be one -- i.e. an Algebra locally \inserted{(fppf)} isomorphic to $M_2(\mathcal O_S)$ -- by
\[
(9)\qquad
\left\{
\begin{array}{rcl}
\unsure{\mathbb G|G}(M)&=&\text{scheme of units of }M\\
S\unsure{\mathbb G|G}(M)&=&\rule{4.8cm}{0.3pt}\quad\text{of determinant }1\\
P\unsure{\mathbb G|G}(M)&=&S\unsure{\mathbb G|G}(M)/\mathrm{Centre}\simeq\unsure{\mathbb G|G}(M)/\unsure{\mathbb G_m|G_m}
\end{array}
\right.
\]
\unsure{Thus|In} the datum \deleted{of $\rho$ defines} \unsure{corresponds} to that of an invertible section of $M$, and
\mspage{257}

\scanpage{259}
the regularity condition means that the subalgebra generated on each fiber is of rank 2 \inserted{(and not reduced to the scalars)} -- which implies that $\rho$ generates \deleted{\illegible{2}}\inserted{\unsure{a|unique} sub}-algebra of degree two
\[
(10)\qquad \mathcal O_S+\mathcal O_S\rho=A_\rho
\]
which is locally a direct \unsure{summand} in $M$ as a module, i.e. with formation \unsure{commuting|commuting} with base change \ldots\ And we shall have
\[
(11)\qquad
\begin{array}{l}
T_\rho=\text{scheme of units of }A_\rho\\
\qquad\text{(case of }G=\unsure{\mathbb G|G}(M)\text{)}
\end{array}
\]
or
\[
(11\text{bis})\qquad
\begin{array}{l}
T_\rho=\text{scheme of units of }A_\rho\\
\qquad\text{of \uline{norm} equal to }1\\
\qquad\text{(case of }G=S\unsure{\mathbb G|G}(M)\text{)}
\end{array}
\]
We shall note that $T_\rho$ is a \uline{torus} \inserted{(necessarily \unsure{maximal})} \unsure{iff} $\rho$ is $\frac{1}{2}$ simple, i.e. $\frac{1}{2}$ simple in each fiber -- which is equivalent, since it is regular, to the fact that in each fiber its two eigenvalues are distinct. For
\[
(12)\qquad \rho=\begin{pmatrix}0&1\\-1&\unsure{0|1}\end{pmatrix}\qquad (\rho^6=1,\ \rho^3=-1)
\]
this is \deleted{\unsure{is}} the case in characteristic $\ne 3$ -- \unsure{thus} \unsure{in the|in the} strict sense $T_\rho$ (\unsure{as} a group scheme over $\operatorname{Spec}\mathbb Z$) is \uline{not} a torus -- it
\mspage{258}

\scanpage{260}
\mspage{431}

becomes \deleted{\unsure{unipotent in characteristic 3}} \inserted{a product of a} unipotent with the center $\unsure{\mathbb G_m|G_m}$ in characteristic 3. Likewise, for
\[
\sigma=\sigma_\infty=\begin{pmatrix}0&-1\\1&0\end{pmatrix},\qquad
\sigma_{\unsure{\infty}}^4=1,\quad\sigma^2=-1,
\tag{13}
\]
$\sigma$ is $\unsure{\frac12}$ simple in all characteristics except in characteristic 2,
where it is unipotent -- and $T_\sigma$ is a
torus except in characteristic 2, where it \deleted{becomes} is
a product of the \inserted{center} $\unsure{\mathbb G_m|G_m}$ with a unipotent.

\margin{Let us return to the general case}
When we take the centralizer in
$S\unsure{\mathbb G|G}(M)$, we find
\[
ST_\rho\subset T_\rho,\qquad
ST_\rho=\operatorname{Ker}\left(
T_\rho
\underset{\inserted{\substack{\text{\unsure{is induced|induced}}\\\text{by the determinant}}}}
{\xrightarrow{\text{norm}}}\unsure{\mathbb G_m|G_m}\right)
\tag{13}
\]

The advantage of the $T_\rho$ over $ST_\rho$ is that
every torsor under $T_\rho=\text{scheme of units of }A_\rho$
is locally trivial in the Zariski sense, whereas
this is no longer necessarily so for
a torsor under $ST_\rho$. Thus, when
we set out to solve the equation in $\alpha$
\[
[\alpha,\rho]=\xi,
\]
for given $\xi$ satisfying the conditions of the
proposition, this is possible (Zariski locally) in
$\unsure{\mathbb G|G}(M)$, but perhaps not in $S\unsure{\mathbb G|G}(M)$.

\mspage{259}

\scanpage{261}

Let us return to the special case
\[
\rho=\begin{pmatrix}0&1\\-1&1\end{pmatrix},\qquad A=\mathbb Z
\]
then
\[
\begin{gathered}
A_\rho=\text{ring of integers of the field}\\
\mathbb Q(\sqrt[3]{-1}),\quad \rho\text{ corresponds to}\\
\exp(2i\pi/6)=\frac{1+i\sqrt3}{2}
\end{gathered}
\tag{14}
\]

When we give \unsure{ourselves} $\xi$, \deleted{the conditions of the}
\[
\xi=\begin{pmatrix}A&B\\C&D\end{pmatrix},\quad
\text{whence }\xi\rho=
\begin{pmatrix}-B&A+B\\-D&C+D\end{pmatrix}
\tag{\deleted{\unsure{15}}}
\]
the conditions \deleted{\unsure{of}} b), c) of the proposition
\deleted{become}
\[
\det\xi=1,\quad \operatorname{Tr}(\xi\rho)=\operatorname{Tr}\rho
\]
\deleted{(16)} become
\[
AD-BC=1\qquad C+D-B=1
\tag{16}
\]
-- quite familiar conditions! They imply
the possibility \deleted{of \unsure{finding}} (over an arbitrary scheme)
of finding Zariski locally an $\alpha\in\Gamma\mathrm{Gl}(2)$
such that \deleted{\unsure{satisfies}}
\[
\alpha\rho(\alpha)^{-1}=[\alpha,\rho]=\xi,
\tag{\unsure{*}}
\]
\margin{at least if we suppose that the condition \unsure{of regularity} is not \unsure{violated} on any fiber\ldots}
and over $\mathbb Z$ it is even possible globally,
for the obstruction is an \deleted{\unsure{torsor under}}
\inserted{invertible $A_\rho$-Module, now}
$A_\rho^{\deleted{*}}$ is principal \unsure{unless I am mistaken} i.e. \deleted{\unsure{Pic of}}
\[
\operatorname{Pic}(A_\rho)=1
\tag{17}
\]

\mspage{260}

\scanpage{262}
\mspage{432}

But if $\alpha$ is a solution, we shall have a priori
\[
\det\alpha\in\mathbb Z^*=\{\pm1\}
\tag{18}
\]
and \deleted{every} other solution will be of the form
$\alpha'=\alpha f$, with $f\in A_\rho^*$, whence
\[
\det\alpha'=\det\alpha\cdot\det f=\det\alpha
\]
for
\[
\det f=f\overline f>0\quad\text{hence }\det f=+1,
\]
thus the sign in (18) \deleted{\unsure{is}} does not depend
on the choice of $\alpha$ -- if it is $-1$, one cannot
find a solution of equation (*) in
$\mathrm{Sl}(2,\mathbb Z)$. The sign depends only on $\xi$ i.e. only
on $A,B,C,D$ -- and in fact \deleted{\unsure{if $D$}} only on $B,D$
(which already determine $\xi$, thanks to
formulas (16) \inserted{at least if $D$ \unsure{invertible}}) -- the computations made elsewhere
suggest that this sign is also \deleted{\illegible{1}} equal to the
remainder of $D$ mod 3 -- $1$ if $D\equiv1\ (3)$, $-1$ if
$D\equiv-1\ (3)$.

Returning to the general case, if we have found
a relation
\[
[\alpha,\rho]=\alpha\rho(\alpha)^{-1}=\xi
\qquad\alpha\in\Gamma\unsure{\mathbb G|G}(M),
\tag{19}
\]
in order \deleted{\unsure{that there exist}} to find a solution $\alpha'$ which
is $\in\Gamma S\unsure{\mathbb G|G}(M)$ i.e. $\det\alpha'=1$, we know that
the solutions $\alpha'\in\Gamma\unsure{\mathbb G|G}(M)$ \deleted{\unsure{sts}} are given by
\[
\alpha'=\alpha f\qquad f\in\Gamma\underline{\mathcal O}_S^*,
\]

\mspage{261}

\scanpage{263}

and the condition becomes
\[
\det\alpha'=1\quad\text{i.e.}\quad\det(\alpha f)=1\quad\text{i.e.}
\]
\[
\det(f^{-1})=\det\alpha
\]
thus the necessary and sufficient condition for the existence
of a solution is that the section $\det\alpha$ of
$\mathcal O_S^*$ be in the image of the norm
homomorphism
\[
\unsure{N|\mathcal N}_{A_\rho/A}:A_\rho^*\longrightarrow A^*,
\tag{20}
\]
\margin{N.B. $S=\operatorname{Spec}A$}
i.e. the obstruction is in the group
\[
A^*/\unsure{N|\mathcal N}_{A_\rho/A}A_\rho^*
\tag{21}
\]
which is a particularly well-known invariant
in class field theory!

\margin{\uline{N.B.} $N(c\rho+d)=\unsure{c^2+cd+d^2}$}
Unless I am mistaken, it follows from class field
theory that for $\rho$ given by (12),
and for $A=\underbrace{\mathbb Z_\ell}_{\inserted{\text{integers }\ell\text{-adic}}}$ ($\ell$ a prime number), this
quotient group is zero in all cases,
except for $\deleted{\unsure{\rho|\ell}}\ \ell=3$, where it is isomorphic to
$\mathbb Z/2\mathbb Z$. \deleted{\illegible{1} Let us note}

\margin{arbitrary \unsure{base|Ring}}
Let us note that (still in the special case (12))
if
\[
\alpha=\begin{pmatrix}a&b\\c&d\end{pmatrix}
\tag{22}
\]
immediate computations already made show that

\mspage{262}